% I. Литературен преглед. Задачата на този раздел е да покаже развитието на
% предметната област, включително технологиите, които вече не се използват.
\section{Анализ на предметната област и литературен преглед}
\label{sec:literature}

\subsection{Развитие на динамичното уеб съдържание}

Първото поколение уеб приложения генерира целия документ на сървъра и го изпраща наново
при всяко действие на потребителя. Динамиката е на ниво страница: състоянието живее в
сървъра, браузърът е програма за изобразяване. Технологиите от този период, сред които
Active Server Pages и генерирането на разметка чрез шаблони от страна на сървъра, решават
задачата за променливо съдържание, но не и задачата за променливо съдържание \term{без
презареждане}. Всяка промяна струва един пълен обмен и едно ново построяване на документа.

Второто поколение измества част от състоянието в браузъра. Скриптовият език на
клиента~\cite{ecma262} получава достъп до построеното дърво на документа и може да го
променя след като страницата е заредена. Спецификацията на този достъп днес е
стандартизирана като обектен модел на документа~\cite{dom_ls}, а разметката и вградените
в нея интерфейси като жив стандарт~\cite{html_ls}. Едновременно с това се появява
асинхронната заявка от скрипт, която по-късно се формализира в общ модел за
извличане~\cite{fetch_ls}. От този момент нататък страницата може да поиска данни, без да
се разрушава.

\subsection{Изместени технологии и защо са изместени}

Част от терминологията в учебното съдържание на дисциплината идва от периода, в който
разширяването на браузъра ставаше чрез двоични компоненти. Технологията ActiveX позволява
изпълнението на native код в контекста на страницата и е налична само в един браузър на
една операционна система. Тя решава реален проблем за времето си, защото тогавашният
браузър няма нито графика извън растерни изображения, нито мрежов достъп извън заявката
за документ. Изместена е по три причини: обвързаност с една платформа, модел на
сигурността, който дава на страницата правата на потребителя, и постепенното покриване на
същите нужди от самите стандарти. Аналогично, филтрите и преходите, реализирани като
собствено разширение на един производител, днес са част от каскадните стилови
спецификации и работят във всеки браузър.

Този проект третира изброените технологии \term{исторически}. Всяка от темите на
дисциплината е реализирана със стандартизирания си наследник, а съответствието между
двете е представено в Раздел~\ref{sec:alternatives}. \TODO{списък на конкретните теми от
учебната програма с точното им позоваване, след сверка със силабуса на дисциплината}

\subsection{Двупосочен обмен между браузър и сървър}

Периодичното запитване (\term{polling}) е най-простият начин да се получат нови данни: с
всяка заявка клиентът пита дали има промяна. Цената е компромис между излишни заявки при
кратък интервал и закъснение при дълъг. Задържаното запитване (\term{long polling})
намалява едното за сметка на сложност в сървъра, но остава в рамките на модела „заявка,
отговор“ на HTTP~\cite{rfc9110}.

Двата съвременни механизма за поток от сървъра към клиента са събитията, изпращани от
сървъра (част от разметковия стандарт~\cite{html_ls}), и протоколът
WebSocket~\cite{rfc6455}. Първият остава в рамките на HTTP и е еднопосочен. Вторият е
отделен протокол, който започва като HTTP заявка за надграждане и след успешно
ръкостискане превръща вече установената връзка в двупосочен канал за съобщения.
Сравнението между двата е основната проектна алтернатива на този проект и е разгледано в
Раздел~\ref{sec:alternatives}.

\subsection{Формати за структуриран обмен}

За обмена на структурирани данни в мрежата съществуват две широко приети представяния.
Разширяемият език за разметка~\cite{xml10} носи схема, пространства от имена и утвърден
инструментариум за преобразуване и проверка. Нотацията за обекти на
JavaScript~\cite{rfc8259} е по-кратка, съответства пряко на структурите от данни на
клиентския език и се разчита от браузъра без допълнителна библиотека. Проектът реализира
\term{и двата} формата за един и същ поток от показатели, за да може разликата между тях
да бъде измерена, а не описана по памет. \TODO{показателите, по които двата формата ще
бъдат сравнени: размер на съобщението, време за разбор, обем на кода за обработка}
